% ===== PRÉAMBULE CANONIQUE — à coller VERBATIM en tête de CHAQUE .tex =====
\documentclass[11pt,a4paper]{article}
\usepackage[utf8]{inputenc}\usepackage[T1]{fontenc}\usepackage{lmodern}
\usepackage[french]{babel}
\usepackage{amsmath,amssymb,mathtools}\usepackage{icomma}
\usepackage{siunitx}\sisetup{locale=FR}  % \qty{}{}, \si{}, \ang{}
\usepackage{xcolor}\usepackage{colortbl}
\usepackage{tikz}\usepackage{pgfplots}\pgfplotsset{compat=1.18}
\usepackage[siunitx,europeanresistors]{circuitikz} % circuits (élec.)
\usepackage[most]{tcolorbox}\usepackage{enumitem}\usepackage{array,booktabs}
\usepackage[a4paper,margin=2.2cm]{geometry}
\usetikzlibrary{arrows.meta,calc,angles,quotes,positioning,patterns,%
  backgrounds,shapes.geometric,decorations.pathmorphing,decorations.markings,intersections}
\definecolor{primary}{RGB}{222,49,99}\definecolor{primarydark}{RGB}{150,22,55}
\definecolor{defcol}{RGB}{0,114,178}\definecolor{methcol}{RGB}{52,92,148}
\definecolor{excol}{RGB}{0,135,90}\definecolor{attncol}{RGB}{200,40,55}
\tcbset{boxbase/.style={enhanced,breakable,boxrule=0.9pt,arc=2.5pt,
  left=3mm,right=3mm,top=2.3mm,bottom=2.3mm,fonttitle=\bfseries,coltitle=white}}
% --- boîtes TUTORIEL (noms EXACTS, ne pas renommer) ---
\newtcolorbox{cle}[1][]{boxbase,colback=defcol!6,colframe=defcol,
  title={\textsc{À retenir}\ifx\relax#1\relax\relax\else\textmd{ : \itshape #1}\fi}}
\newtcolorbox{methode}[1][]{boxbase,colback=methcol!7,colframe=methcol,sharp corners,
  title={\textsc{Méthode}\ifx\relax#1\relax\relax\else\textmd{ --- \itshape #1}\fi}}
\newtcolorbox{exemplebox}[1][]{boxbase,colback=excol!5,colframe=excol,
  title={\textsc{Exemple}\ifx\relax#1\relax\relax\else\textmd{ : \itshape #1}\fi}}
\newtcolorbox{piege}[1][]{boxbase,colback=attncol!6,colframe=attncol,
  title={\textsc{Piège}\ifx\relax#1\relax\relax\else\textmd{ : \itshape #1}\fi}}
\newtcolorbox{remarque}[1][]{boxbase,colback=black!4,colframe=black!45,
  title={\textsc{Remarque}\ifx\relax#1\relax\relax\else\textmd{ : \itshape #1}\fi}}
% --- boîtes EXERCICE / CORRIGÉ (noms EXACTS, auto-comptées) ---
\newtcolorbox[auto counter]{exo}[1][]{boxbase,colback=primary!4,colframe=primary,
  title={\textsc{Exercice}~\thetcbcounter\ifx\relax#1\relax\relax\else\textmd{ --- \itshape #1}\fi}}
\newtcolorbox[auto counter]{corrige}[1][]{boxbase,colback=excol!4,colframe=excol,
  title={\textsc{Corrigé}~\thetcbcounter\ifx\relax#1\relax\relax\else\textmd{ --- \itshape #1}\fi}}
\newcommand{\vv}[1]{\vec{#1}}\newcommand{\ds}{\displaystyle}
\newcommand{\R}{\mathbb{R}}\newcommand{\N}{\mathbb{N}}\newcommand{\Z}{\mathbb{Z}}
% =========================================================================
\begin{document}
\sisetup{per-mode=symbol}
% ================= CORRIGÉ FACILE =================

\begin{corrige}[ascenseur à l'arrêt]
À l'arrêt, $a=0$, donc $N=m(g+0)=mg$ :
\[
N=80\times 9{,}8=\qty{784}{\newton}.
\]
La balance affiche le poids normal (soit une masse apparente de $\qty{80}{\kilogram}$).
\end{corrige}

\begin{corrige}[l'ascenseur accélère vers le haut]
$a>0$, donc $N>mg$ :
\[
N=m(g+a)=60\times(10+2)=60\times 12=\qty{720}{\newton}.
\]
On se sent plus lourd ($720>600$).
\end{corrige}

\begin{corrige}[l'ascenseur accélère vers le bas]
$a<0$, donc $N<mg$ :
\[
N=m(g+a)=60\times(10-2)=60\times 8=\qty{480}{\newton}.
\]
On se sent plus léger ($480<600$).
\end{corrige}

\begin{corrige}[le câble casse]
En chute libre, $a=-g$ :
\[
N=m(g-g)=50\times 0=\qty{0}{\newton}.
\]
La balance affiche $0$ : la personne est en \textbf{impesanteur}, elle flotte dans la
cabine. Son poids réel $mg=\qty{490}{\newton}$ existe pourtant toujours.
\end{corrige}

\begin{corrige}[deviner l'état de l'ascenseur]
On compare $N$ à $mg=\qty{700}{\newton}$.
\begin{enumerate}[label=\alph*)]
  \item $N=mg$ : ascenseur \textbf{au repos ou en MRU} ($a=0$).
  \item $N>mg$ : ascenseur qui \textbf{accélère vers le haut} ($a>0$).
  \item $N<mg$ : ascenseur qui \textbf{accélère vers le bas} ($a<0$).
\end{enumerate}
\end{corrige}

\end{document}
